% 全书各部分、章和节标题的统一定义。
\BookPart{第一部分}{在一个事务域中把保证做实}

\BookChapter{01}{先实现一个正常工作的下单服务}{%
  \BookSection{订单、账户、扣款流水与购买操作}
  \BookSection{下单与本地扣款的事务}
  \BookSection{购买流程、支付动作与 HTTP 请求}
  \BookSection{单数据库、单付款路径与最小货币单位}
}

\BookChapter{02}{一次超时，为什么变成了两次扣款}{%
  \BookSection{制造一次提交后的响应丢失}
  \BookSection{客户端、服务端与数据库的时间线}
  \BookSection{请求未到达、执行回滚与提交后失联}
  \BookSection{业务效果、请求尝试与可观察结果}
}

\BookChapter{03}{系统怎样认出同一次购买}{%
  \BookSection{为同一次业务意图分配稳定身份}
  \BookSection{请求 ID、链路 ID、订单 ID 与幂等键}
  \BookSection{命名空间、租户、动作与键的生成和传播}
  \BookSection{同键异参与两次合法购买}
  \BookSection{第一版 API 规则：有效期与身份假设}
}

\BookChapter{04}{两个请求同时到达，第一版防重失效了}{%
  \BookSection{用并发屏障击穿“先查后写”}
  \BookSection{唯一约束与事务中的仲裁}
  \BookSection{冲突、等待、回滚与隔离级别}
  \BookSection{重试之间插入另一笔合法交易}
  \BookSection{数学幂等、交换律与业务正确性}
}

\BookChapter{05}{扣款已经成功，怎样把结果找回来}{%
  \BookSection{稳定结果与当前状态查询}
  \BookSection{把操作事实、流水与结果一起提交}
  \BookSection{重放原结果，保留每次响应的链路信息}
  \BookSection{终局失败、临时故障与同键异参}
  \BookSection{投递、执行、提交与生效的次数}
}

\BookChapter{06}{在每一步杀死服务，检验第一版协议}{%
  \BookSection{操作状态机与完整执行历史}
  \BookSection{线性化点：提交与获知结果之间}
  \BookSection{不重复生效、不错误合并与结果一致的证明}
  \BookSection{进展需要哪些可用性与重试条件}
  \BookSection{让错误协议在有界模型中留下反例}
}

\BookPart{第二部分}{走出本地事务，学会处理未知结果}

\BookChapter{07}{接入支付平台，本地事务的保证到哪里结束}{%
  \BookSection{换用外部支付后的提交边界}
  \BookSection{下游的幂等键与结果查询能力}
  \BookSection{查询暂未查到，旧请求却仍可能到达}
  \BookSection{用固定操作身份串起回调、查询与对账}
  \BookSection{不同下游能力下的保证边界}
}

\BookChapter{08}{支付一直停在处理中，谁来接着做}{%
  \BookSection{把恢复任务保存为持久化工作流}
  \BookSection{业务状态、工作者所有权与外部结果的确定性}
  \BookSection{租约、版本与接管：旧工作者恢复之后}
  \BookSection{资源端的 fencing token 检查与写入保护边界}
  \BookSection{Leases、Chubby 与 Redlock 的系统假设}
}

\BookChapter{09}{支付成功以后，可靠地通知下游}{%
  \BookSection{数据库与消息系统之间的双写缝隙}
  \BookSection{将 Outbox 事件与业务结果一起提交}
  \BookSection{发布成功、标记未完成时重启}
  \BookSection{事件身份、发布顺序、重试与积压}
}

\BookChapter{10}{消费者处理完消息，却没来得及确认}{%
  \BookSection{先 ACK、先修改还是先去重}
  \BookSection{在消费者事务中提交去重记录与业务效果}
  \BookSection{重复、乱序、重放与消费者升级}
  \BookSection{消息事务与外部数据库、支付副作用的边界}
  \BookSection{整条链路上的次数语义}
}

\BookPart{第三部分}{复制、遗忘和历史变化之后，承诺是否还在}

\BookChapter{11}{切换到另一个副本，系统还记得做过什么吗}{%
  \BookSection{数据库高可用与复制状态机的两条路径}
  \BookSection{Raft 如何用日志、提交与应用保留操作事实}
  \BookSection{状态机中的客户端会话、操作序列与去重结果}
  \BookSection{重复日志项、重放、快照与日志截断}
  \BookSection{回到 Paxos、Raft 与主备复制的原始资料}
}

\BookChapter{12}{去重记录越来越多，能不能删掉}{%
  \BookSection{有限重试约定与记录保留策略}
  \BookSection{结果回收、业务证据与最小拒绝证据}
  \BookSection{序列号、水位线、会话与 epoch 的成立条件}
  \BookSection{未完成请求、序列空洞与旧会话重连}
  \BookSection{超期以后：拒绝、查询或开启新操作}
}

\BookChapter{13}{换一种业务表达，能少记住一些历史吗}{%
  \BookSection{购买系统中的赋值、增量、集合合并与条件迁移}
  \BookSection{分别检验幂等、交换律、单调性与业务不变量}
  \BookSection{状态型 CRDT 的合并性质与操作型 CRDT 的传输条件}
  \BookSection{标签与履约进度：何时可以少记一些历史}
  \BookSection{扣款、发券与退款为何仍可能需要操作身份}
}

\BookChapter{14}{恢复旧备份，再把流量切到另一个地域}{%
  \BookSection{分片路由与操作归属：同键在多个分片间的受理规则}
  \BookSection{崩溃恢复、无损切换与旧历史回退}
  \BookSection{外部支付账本、事件日志与本地操作证据}
  \BookSection{单归属路由、跨地域仲裁与分区期间的受理条件}
  \BookSection{恢复后的重试入口与对账完成前的行为}
}

\BookChapter{15}{给最终系统写下它能兑现的保证}{%
  \BookSection{不同部署与下游能力下的最终保证}
  \BookSection{退款协议、额度仲裁与跨版本重放}
  \BookSection{业务不变量、进展前提与无法处理的历史}
  \BookSection{把论证、测试、模型检查与恢复演练连起来}
  \BookSection{沿核心保证查找保存来源与验证证据}
}

\BookPart{第四部分}{分布式购买系统的整体设计与工程架构}

\BookChapter{16}{总体架构、服务协作与业务链路}{%
  \BookSection{系统设计条件与全局视图}
  \BookSection{数据所有权与跨域业务协议}
  \BookSection{支付、事件与履约关键链路}
  \BookSection{RabbitMQ 消息子系统}
  \BookSection{Kafka 日志与消费子系统}
}

\BookChapter{17}{分区数据、复制与协调控制面}{%
  \BookSection{逻辑分区与业务事实归属}
  \BookSection{完整分区的 Raft 复制与 Chain Replication}
  \BookSection{配置、会话与执行资格}
  \BookSection{完整业务分区的在线迁移}
  \BookSection{TiKV 的 Region、Peer 与事务边界}
}

\BookChapter{18}{Redis 节点、缓存与集群设计}{%
  \BookSection{请求执行与查询缓存}
  \BookSection{RDB 快照与恢复加载}
  \BookSection{AOF、重写与启动选择}
  \BookSection{复制、等待与 Sentinel}
  \BookSection{Cluster 的槽、拓扑与迁移内部}
}

\BookChapter{19}{分布式存储与跨分区事务}{%
  \BookSection{块、副本与元数据}
  \BookSection{完整图像与发布结构}
  \BookSection{跨分区资源与原子决定}
  \BookSection{Spanner 与分布式 SQL 事务架构}
  \BookSection{FaRM 的内存、RDMA 与事务路径}
}

\BookChapter{20}{批处理、任务、Actor 与函数执行}{%
  \BookSection{固定输入与批量对账数据流}
  \BookSection{重执行、推测与输出发布}
  \BookSection{Ray 的任务、对象与调度}
  \BookSection{Actor 状态与业务恢复}
  \BookSection{Lambda 事件调用与执行环境}
}

\BookChapter{21}{多地域部署、容量运维与前沿架构}{%
  \BookSection{地域部署与隔离恢复}
  \BookSection{容量、热点与恢复预算}
  \BookSection{可观测性、运维控制与恢复职责}
  \BookSection{成熟工业实现}
  \BookSection{近期前沿架构}
}
